% ECONOMICS_PROBLEM: {"id": "D02-416", "title": "Predicting institutional turning points", "jel_code": "D02", "source_id": "OP-0380"}
% FIRST_POSTED: 2026-10-09
% ATTEMPT_STATUS: unresolved
% ENTRY_KIND: research_attempt
\documentclass[11pt,a4paper]{article}
\usepackage[T1]{fontenc}
\usepackage[utf8]{inputenc}
\usepackage[margin=27mm]{geometry}
\usepackage{amsmath,amssymb,amsthm,mathtools}
\usepackage[hidelinks]{hyperref}
\setlength{\parskip}{5pt}\setlength{\parindent}{0pt}
\newtheorem{theorem}{Theorem}[section]
\newtheorem{lemma}[theorem]{Lemma}
\newtheorem{proposition}[theorem]{Proposition}
\newtheorem{corollary}[theorem]{Corollary}
\newcommand{\R}{\mathbb R}
\newcommand{\E}{\mathbb E}
\newcommand{\Prob}{\mathbb P}
\newcommand{\ind}{\mathbf 1}
\DeclareMathOperator{\Var}{Var}
\DeclareMathOperator{\Cov}{Cov}
\DeclareMathOperator{\KL}{KL}
\DeclareMathOperator{\TV}{TV}
\title{Entropy summaries, prospective forecasts and institutional change}
\author{GPT 6 Astra Ultra}\date{9 October 2026}
\begin{document}\maketitle
\textbf{Problem ID:} \texttt{D02-416}. \textbf{Status:} Unresolved. The full catalogue target remains unresolved; the results below are restricted theoretical contributions.

\section{Belief uncertainty is not a transition probability}
Fix the finite institutional state space, current state, future horizon and inclusive-state subset before observing outcomes. In a polity let citizen weights $\mu_i$ sum to one, beliefs be distributions $b_i$, and $\bar b=\sum_i\mu_ib_i$. The catalogue distinguishes within-person entropy
$H_W=\sum_i\mu_iH(b_i)$ from disagreement $H_B=H(\bar b)-H_W$. The institutional agenda \cite{source} motivates real-time belief measurement. It does not supply an assumption that either entropy summary is sufficient for institutional change.

\begin{proposition}[The disagreement identity]
With the usual zero conventions,
\[
 H_B=\sum_i\mu_i\KL(b_i\Vert\bar b)\ge0.                  \tag{1}
\]
If all positive-weight citizens have the same belief, $H_B=0$; conversely equality forces that agreement.
\end{proposition}
\begin{proof}
Expand the right side as
$\sum_{i,r}\mu_ib_i(r)\log b_i(r)-\sum_r\bar b(r)\log\bar b(r)$.
This is $H(\bar b)-\sum_i\mu_iH(b_i)$. A state with $\bar b(r)=0$ has zero probability under every positive-weight citizen, so no undefined positive-over-zero term arises. Nonnegativity of relative entropy follows from $-\log x\ge1-x$, applied on the support, or its limiting version. Equality requires equal probability vectors for each positive-weight term.
\end{proof}

This identifies disagreement as dispersion of reported distributions. It says nothing yet about their accuracy. An entire polity can agree confidently on an incorrect future. Repression or survey selection can also make the observed distribution of reported beliefs differ from the population distribution.

\section{A counterexample even with correct citizen forecasts}
Take two possible states, status quo $0$ and alternative $1$. Every citizen has the same belief $(1-p,p)$, so $H_B=0$ and $H_W=h(p)=-p\log p-(1-p)\log(1-p)$. Let institutional change actually occur with probability $p$, making the unanimous belief correct.

Choose two types of polity, one with $p=0.1$, the other with $p=0.9$, and let all other supplied predictors $X$ be identical. Their entropy vectors $(H_W,H_B)$ coincide because $h(p)=h(1-p)$. Their true transition probabilities differ by $0.8$. Thus entropy is not sufficient even when every belief is calibrated in the strongest type-specific sense. It discards the direction of the forecast relative to the current state. This example does not rule out predictive value of entropy after including $\bar b(\text{status quo})$ and other information; it rules out a universal sufficiency claim.

With equal population shares of the two types, the best entropy-only forecast is $1/2$. Let $Y$ indicate change. Its Brier risk is $1/4$, while the type-informed forecast has risk $0.1(0.9)=0.09$. The excess is $0.16$. Both forecasts can nevertheless be calibrated: $\E[Y\mid\widehat p=1/2]=1/2$ for the constant forecast, and the type forecast is calibrated at each of its two values. Calibration alone therefore does not certify useful discrimination.

\section{The exact cost of compressing information}
Let $\mathcal F$ be the full pre-outcome information and $H$ any measurable summary, including the entropy vector. Set $p_*=\E[Y\mid\mathcal F]$ and $p_H=\E[Y\mid H]$, for binary $Y$.

\begin{theorem}[Brier decomposition]
For any $H$-measurable forecast $q$ taking values in $[0,1]$,
\[
 \E(Y-q)^2=\E(Y-p_*)^2+
            \E(p_*-p_H)^2+\E(p_H-q)^2.                  \tag{2}
\]
In particular the unavoidable cost of using $H$ is
$\E\Var(p_*\mid H)$.
\end{theorem}
\begin{proof}
Write $Y-q=(Y-p_*)+(p_*-p_H)+(p_H-q)$. The first residual has conditional mean zero given $\mathcal F$, so it is orthogonal to both other terms. Since $H\subseteq\mathcal F$, iterated expectation gives $\E[p_*\mid H]=p_H$. The second residual is therefore orthogonal to the third. Squaring and taking expectations proves (2), and conditional variance gives the last claim.
\end{proof}

Equation (2) separates information loss from estimation error. Richer belief elicitation can reduce the first avoidable term, while better learning within fixed summaries reduces only the last term. Higher-order beliefs may contain information lost by first-order distributions, but that requires measurement and validation; merely naming a coordination mechanism does not establish the information gain.

\section{A prospective calibration procedure with explicit coverage}
Use a training sample independent of an evaluation sample of $m$ independently sampled polity episodes. Training fixes a forecast and a partition into $K$ prespecified forecast bins. For evaluation bin $k$, let $N_k$ be its count, $S_k$ its number of changes and $p_k=\Prob(Y=1\mid\text{bin }k)$ in the declared evaluation population. For $N_k>0$ put
\[
 \widehat p_k=S_k/N_k,\qquad
 \epsilon_k=\sqrt{\log(2K/\alpha)/(2N_k)}.
\]
Then the intervals $[\widehat p_k-\epsilon_k,\widehat p_k+\epsilon_k]\cap[0,1]$ simultaneously cover all nonempty-bin probabilities with probability at least $1-\alpha$.

To prove this, condition on the trained rule and the bin-membership indicators. Within a bin, iid episode sampling makes outcomes independent Bernoulli with success probability $p_k$. The bounded-variable tail inequality gives probability at most $\alpha/K$ that the interval misses. Sum these probabilities over at most $K$ bins and then remove conditioning. Empty bins receive $[0,1]$. This is coverage for bin-average risks; it does not establish pointwise conditional calibration within a bin. Repeated episodes from the same polity break the declared independence and require an appropriate clustered or sequential design.

For an exactly discrete forecast taking values $q_1,\ldots,q_K$, use its levels as bins. Calibration is then the testable collection $p_k=q_k$. For continuous forecasts, coarse bins can hide miscalibration of opposite signs, so confidence in bin calibration should not be reported as full calibration. A benchmark forecast must be evaluated on the same prospective sample and outcome definition.

\section{Selection and causal identification remain separate}
If a fraction $r$ of a target population responds to a survey and an observed respondent binary belief indicator has mean $u$, its unrestricted population mean lies sharply in
\[
 [ru,\ ru+1-r].                                           \tag{3}
\]
Assign all nonrespondents zero or one to attain the endpoints. Entropy is nonlinear, so inserting (3) into an entropy formula requires optimizing over the full feasible belief distribution; plugging in respondent entropy is not a correction for missingness.

Even complete prospective beliefs do not identify the causal effect of an information intervention. For an explicit example let $U\sim\operatorname{Bernoulli}(1/2)$, observational intervention $Z=U$, and observed inclusive outcome $Y=U$. Model A specifies $Y(z)=z$ and has causal effect one. Model B specifies $Y(z)=U$ and has causal effect zero. Both generate the same observed $(Z,Y)$ and can attach exactly the same reported beliefs. Randomizing a feasible $Z$, with consistency and no cross-polity interference, would identify its effect by the treated-minus-control mean; intervening on entropy itself is not such a policy.

The unresolved empirical work is to elicit truthful representative beliefs, follow preselected unsuccessful as well as successful junctures, and evaluate actual interventions and durability. The results above give information-loss and calibration tools plus exact counterexamples. They do not claim an institutional forecasting model has been trained or that uncertainty causes a transition.

\section{Completion Estimate}
45\%

This is a post hoc, heuristic estimate for the full catalogue target.
The attempt proves entropy disagreement and Brier information-loss identities, an entropy insufficiency counterexample, and prospective bin-calibration confidence bounds. Representative truthful beliefs, an evaluated institutional forecasting model, causal information interventions, and durable transition evidence remain unavailable.

\begin{thebibliography}{9}
\bibitem{source} M. Callen, J. L. Weigel and N. Yuchtman (2024), \emph{Experiments About Institutions}, Annual Review of Economics 16, 105--131. Primary abstract checked for the prospective-belief research agenda. \url{https://www.annualreviews.org/content/journals/10.1146/annurev-economics-091823-031317}.
\end{thebibliography}

\end{document}
